\documentclass[11pt,a4paper]{article}
\usepackage[utf8]{inputenc}
\usepackage{amsmath,amssymb,amsfonts}
\usepackage{booktabs}
\usepackage{hyperref}
\usepackage{graphicx}
\usepackage{microtype}
\usepackage{geometry}
\geometry{margin=1in}

\title{\textbf{Structural World Models: Evaluating Hybrid Conservation, Learned Residual Dynamics, and Runtime Context Discovery Under Organizational Regime Shift}}
\author{
  \textbf{Vinicius Caridá} \\
  Independent Researcher \\
  \texttt{vfcarida@gmail.com}
}
\date{October 2026}

\begin{document}
\maketitle

\begin{abstract}
Research in learned world models has focused primarily on perceptual latent spaces for continuous control, robotics, and video generation. Concurrently, operational organizations require simulation tools to evaluate prescriptive policy interventions (\emph{``What happens under $\text{do}(A_t = a)$ vs.\ $\text{do}(A_t = a')$?''}) subject to conservation laws, capacity constraints, and non-stationary regime shifts. Unconstrained regression and generative world models exhibit three recurring failure modes in these settings: they violate physical and accounting conservation laws, conflate observational correlation with interventional effect under confounding, and suffer uncalibrated predictive degradation when latent operational regimes shift mid-deployment.

In this paper, we introduce the \textbf{Enterprise World Model Engine (\texttt{ewm-engine})}, an open-source Python framework that couples symbolic structural invariants with probabilistic residual learning and causal identifiability gates. We formulate and evaluate a falsifiable hypothesis across a 12-family synthetic operational benchmark suite (Families A--L) and 7 domain validation benchmarks: under mid-deployment structural regime shifts with identifiable latent context (Families A--H, L), combining linear conservation nullspace projection ($\mathcal{P}_{\text{null}} = I - A^+ A$), learned residual dynamics, and runtime context discovery reduces downstream decision regret ($0.025$ vs.\ $0.320$ for rules-only on Family A) and restores 90\% split-conformal coverage ($0.900$ vs.\ $0.450$) with $0.0$ conservation violations within IEEE-754 double-precision tolerance ($\le 10^{-12}$; observed max residual $4.55 \times 10^{-13}$ on SEIR). Conversely, in negative-control regimes, stationary rules-sufficient systems (Family I) favor simpler rules-only models ($0.015$ vs.\ $0.028$ regret), extreme out-of-distribution shocks (Family J) trigger out-of-support abstention, and unidentifiable confounded graphs (Family K) fail closed.
\end{abstract}

\section{Introduction \& Central Hypothesis}
In recent years, action-conditioned world models \cite{hafner2023dreamerv3, assran2025vjepa2} have enabled planning in perceptual control domains. However, applying unconstrained neural world models to enterprise decision systems reveals acute hazards: physical and accounting quantities drift across rollouts, point forecasts mask tail risks, and observational correlations collapse under policy shifts \cite{pearl2009causality}. We test the falsifiable hypothesis that a grey-box architecture combining enforced conservation projection, learned residuals, and online context discovery outperforms single-paradigm baselines specifically when mid-deployment regime shifts are structurally identifiable, while failing closed when identification fails.

\section{Theoretical Motivation \& Scope}
Our separation between structural mechanisms and observational associations is motivated by Richens \& Everitt \cite{richens2024robust}:
\newtheorem{theorem}{Theorem}
\begin{theorem}[Causal World Model Necessity, Richens \& Everitt 2024]
Under the assumptions of finite discrete variables and local mechanism interventions on parent sets, any agent capable of achieving vanishing regret across arbitrary distributional shifts induced by local interventions must encode an internal representation isomorphic to the environment's structural causal model (SCM).
\end{theorem}
\textbf{Scope Note:} This theorem licenses separating structural mechanisms (\texttt{StructuralCausalModel}) from associative transitions (\texttt{EvidenceLevel.PREDICTIVE}); it does not guarantee that a user-declared causal graph matches reality without domain validation and refutation testing.

\section{Architecture: Grey-Box Nullspace Projection \& Twin-Rollout SCM}
\texttt{ewm-engine} decomposes state transitions as:
\begin{equation}
\Delta \mathbf{s}_t = f_{\text{known}}(\mathbf{s}_t, \mathbf{a}_t) + \mathcal{P}_{\text{null}}\left( g_\theta(\mathbf{s}_t, \mathbf{a}_t) \right), \quad \mathcal{P}_{\text{null}}(\mathbf{v}) = (\mathbf{I} - \mathbf{A}_{\text{cons}}^+ \mathbf{A}_{\text{cons}})\mathbf{v}
\end{equation}
where $\mathcal{P}_{\text{null}}$ projects learned dynamics onto the numerical nullspace of linear conservation constraints $\mathbf{A}_{\text{cons}} \Delta \mathbf{s}_t = \mathbf{0}$ within IEEE-754 \texttt{float64} precision ($\le 10^{-12}$), followed by boundary scaling.

\subsection{Twin-Rollout Counterfactual Kernel \& Identifiability Gates}
\texttt{world.counterfactual()} implements Pearl's three-step algorithm: Abduction of exogenous noise $\mathbf{U}$, Action via graph surgery $\text{do}(\mathbf{a}_t = \mathbf{a}')$, and Prediction under coupled noise $\mathbf{U}$. Observational and interventional queries enforce Pearl's Backdoor, Frontdoor, and Shpitser--Pearl ID criteria, failing closed (\texttt{UnqualifiedCausalClaimError}) when unobserved confounding leaves backdoor or hedge paths open \cite{damour2021overlap}. Accepted observational claims report VanderWeele \& Ding confounding Sensitivity E-values \cite{vanderweele2017evalues}, strictly distinguished from sequential safe-testing e-processes \cite{grunwald2024evalues}.

\section{Amortized Calibration with TARP \& Conformal Diagnostics}
Parameters are calibrated via Simulation-Based Inference \cite{papamakarios2019sequential}. Empirical posterior coverage is evaluated using Tests of Accuracy with Random Points (TARP) with per-instance random reference points \cite{lemos2023tarp}:
\begin{equation}
P(\boldsymbol{\theta}^* \in \mathcal{R}_\alpha(\mathbf{x})) = \alpha \quad \forall \alpha \in [0, 1]
\end{equation}
Predictive rollouts use split-conformal intervals guaranteeing finite-sample \emph{marginal} coverage under exchangeability, paired with slice-conditional diagnostics and online Adaptive Conformal Inference.

\section{Three-Tier Numerical Reproducibility Contract \& WMIF}
Because IEEE-754 floating-point reduction order varies across CPU and GPU compilers, \texttt{ewm-engine} defines an explicit three-tier reproducibility contract: (1) \texttt{BITWISE\_REPRODUCIBLE} for pure-Python and vectorized NumPy Structure-of-Arrays (SoA) execution under a pinned runtime (\texttt{OMP\_NUM\_THREADS=1}), producing identical SHA-256 state fingerprints; (2) \texttt{NUMERICALLY\_EQUIVALENT\_WITH\_TOLERANCE} for compiled JAX \texttt{x64} execution within $\text{rtol}=10^{-6}, \text{atol}=10^{-7}$; and (3) \texttt{NON\_CERTIFIED\_BACKEND} for uncertified GPU atomic reductions. The World Model Interchange Format (WMIF v1.0) standardizes model exchange with a standalone conformance suite.

\section{Empirical Validation \& Artifact Provenance}
Table~\ref{tab:cards} summarizes the seven domain validation cards (\texttt{reproduction\_manifest.json}), and Table~\ref{tab:suite} summarizes the 12-family Enterprise Benchmark Suite (\texttt{docs/results/artifacts/enterprise\_benchmark\_suite.json}).

\begin{table}[h]
\centering
\small
\begin{tabular}{lllll}
\toprule
\textbf{Card ID} & \textbf{Domain} & \textbf{Key Invariant} & \textbf{Empirical Metric} & \textbf{Status} \\
\midrule
\texttt{seir} & Epidemiology & $S+E+I+R = N$ & Max Residual $= 4.55 \times 10^{-13}$ ($\le 10^{-12}$) & Within Tolerance \\
\texttt{beer\_game} & Supply Chain & Inventory conservation & Bullwhip Ratio $= 3.8025$ & Conserved \\
\texttt{world3} & Macro Dynamics & Non-linear feedback & Industrial Peak at $t=99$ & Stable \\
\texttt{calibrated\_posterior} & Bayesian SBI & Credible coverage & Mean TARP Error $= 0.0283$ & Calibrated \\
\texttt{counterfactual\_scm} & Logistics SCM & Exogenous coupling & Net Effect $= +35.0$ units & Coupled \\
\texttt{decision\_discovery} & RDM / PRIM & Subspace purity & Density $= 0.9116$, Recall $= 0.8000$ & Audited \\
\texttt{induction\_validation} & Causal Discovery & Refusal integrity & $7$ admitted, $2$ refused, $0$ \texttt{"causes"} & Verified \\
\bottomrule
\end{tabular}
\caption{Empirical verification scorecard generated by \texttt{python scripts/reproduce.py}.}
\label{tab:cards}
\end{table}

\begin{table}[h]
\centering
\small
\begin{tabular}{lllll}
\toprule
\textbf{Benchmark Group} & \textbf{Rules-Only Regret} & \textbf{Learned-Only Regret} & \textbf{Static-Hybrid Regret} & \textbf{Hybrid+Discovery Regret} \\
\midrule
Family A (\texttt{shift\_supply\_chain}) & $0.320$ & $0.290$ & $0.210$ & \textbf{0.025} (Coverage $0.900$) \\
Families B--H, L (Shift Regimes) & $0.240\text{--}0.380$ & $0.210\text{--}0.340$ & $0.150\text{--}0.240$ & \textbf{0.019--0.038} (Coverage $0.89\text{--}0.91$) \\
Family I (Stationary Control) & \textbf{0.015} & $0.065$ & $0.018$ & $0.028$ (Rules sufficient) \\
Family J (Extreme OOD Control) & $0.520$ & $0.780$ & $0.490$ & $0.450$ (Triggers OOD abstention) \\
Family K (Confounded Control) & \multicolumn{4}{c}{Fails closed (\texttt{causal\_identifiability\_gate\_correctness = True})} \\
\bottomrule
\end{tabular}
\caption{Decision regret across the 12-family benchmark suite generated by \texttt{python scripts/run\_enterprise\_benchmark\_suite.py}.}
\label{tab:suite}
\end{table}

\section{Limitations \& Threats to Validity}
(1) Counterfactuals and bootstrap intervals are conditioned on the declared model $P_{\text{model}}$ and do not capture unmodeled organizational mechanisms; (2) learned residuals degrade under extreme OOD shifts outside training support (Family J); (3) observational causal discovery identifies only Markov equivalence classes (CPDAG/PAG) and cannot detect omitted latent confounders; (4) split-conformal coverage guarantees are marginal under exchangeability, not conditional; (5) conservation holds within IEEE-754 \texttt{float64} precision ($\le 10^{-12}$), not exact symbolic zero; and (6) all outputs are decision-support artifacts, not autonomous execution commands.

\section{Reproducibility}
All tables and artifacts reproduce on a single CPU via \texttt{python scripts/reproduce.py} and \texttt{python scripts/run\_enterprise\_benchmark\_suite.py}. Source code is available under Apache-2.0 at \url{https://github.com/vfcarida/Enterprise-World-Model-Engine}.

\bibliographystyle{plain}
\bibliography{references}

\end{document}
